\documentclass[12pt]{article}
\usepackage[margin=1in]{geometry}
\usepackage{fontspec}
\usepackage{unicode-math}
\setmainfont[
  Path=manuscript/fonts/,
  BoldFont=DejaVuSerif-Bold.ttf,
  ItalicFont=DejaVuSerif-Italic.ttf,
  BoldItalicFont=DejaVuSerif-BoldItalic.ttf
]{DejaVuSerif.ttf}
\setmathfont[Path=manuscript/fonts/]{DejaVuMathTeXGyre.ttf}
\usepackage{microtype}
\usepackage{lineno}
\input{manuscript/data/cross_document_refs.tex}
\setlength{\emergencystretch}{2em}
\clubpenalty=10000
\widowpenalty=10000
\displaywidowpenalty=10000
\usepackage{amsmath}
\usepackage{graphicx}
\usepackage{booktabs}
\usepackage{array}
\setlength{\tabcolsep}{3pt}
\newcolumntype{L}[1]{>{\raggedright\arraybackslash}p{#1}}
\usepackage{longtable}
\usepackage{enumitem}
\usepackage{titling}
\setlength{\droptitle}{-2em}
\posttitle{\par\end{center}}
\predate{}\postdate{}
\usepackage{titlesec}
\titleformat{\section}{\normalfont\Large\bfseries\raggedright}{}{0pt}{}
\titleformat{\subsection}{\normalfont\large\bfseries\raggedright}{}{0pt}{}
\titleformat{\subsubsection}{\normalfont\normalsize\bfseries\raggedright}{}{0pt}{}
\usepackage{needspace}
\usepackage[section]{placeins}
\usepackage{hyperref}
\usepackage{xurl}
\usepackage{xcolor}
\usepackage{float}
\renewcommand{\topfraction}{0.90}
\renewcommand{\bottomfraction}{0.80}
\renewcommand{\textfraction}{0.08}
\renewcommand{\floatpagefraction}{0.75}
\setcounter{topnumber}{3}
\setcounter{totalnumber}{4}
\usepackage{caption}
\setlength{\parskip}{0.35em}
\setlength{\parindent}{0pt}
\hypersetup{colorlinks=true,linkcolor=blue,citecolor=blue,urlcolor=blue}

\title{TRiX: Task-conditioned execution governance for robotic disassembly}
\author{Stavan Dholakia\textsuperscript{1}, Shivani Shukla\textsuperscript{2*}\\
Abhishek Singh\textsuperscript{3}, Aditya Gazta\textsuperscript{4}}
\date{}

\begin{document}
\begingroup
\setlength{\parskip}{0pt}
\maketitle
\endgroup

{\footnotesize
\noindent\textsuperscript{1}Cloud and AI, Microsoft, One Microsoft Way, Redmond, WA 98052, USA.\par
\noindent\textsuperscript{2}Analytics and Information Systems, University of San Francisco, 2130 Fulton St, San Francisco, CA 94117, USA.\par
\noindent\textsuperscript{3}GCP, Google, 15303 Dallas Pkwy, Addison, TX 75001, USA.\par
\noindent\textsuperscript{4}Robotics, Amazon, 13455 Noel Rd, Dallas, TX 75240, USA.\par
\noindent\textsuperscript{*}Corresponding author: sgshukla@usfca.edu.\par
}

\begin{abstract}
Robots disassembling electronics increasingly act on commands
proposed by learned policies and vision-language models. A plausible command
can fracture a circuit board, strip a thread or break a latch. Here we
show that the value of a runtime governor depends on whether the proposing
policy was trained with it. We present TRiX, which allows, corrects
or rejects each command against a set conditioned on task state, and test it
in six simulated tasks, on recorded model commands and on a robot arm. In circuit-board extraction, one restriction gives 0\%
safe completion when added after training and 100\% when present during
training. In crank extraction, training cannot recover the authority a fixed
restriction removes: 0\% against 49.4\%. Under snap-fit conditions outside
the training range, safe completion falls from 96.5\% to 59.9\% although every
executed command stays within its bounds. Across 273 recorded model decisions,
governance prevents all 31 observed constraint-violation cases. Governors must
be evaluated with the policies they govern.
\end{abstract}

\noindent\textbf{Keywords:} robotic disassembly; execution governance; action
projection; robot learning; vision-language models
\linenumbers

\section{Introduction}
\label{sec:intro}
Electronic waste combines heterogeneous products with manipulation failures
that can permanently damage valuable or hazardous components
\cite{ref1,ref2,ref3,ref4,ref5,ref6}. A robot that disassembles such products
must decide, at every step, what force or motion to apply to a part whose
condition keeps changing. Vision-language models (VLMs), language models and
learned policies now let a robot propose a wide range of such actions from
images, instructions and experience \cite{ref8,ref9,ref10,ref11,ref12,ref13}.
Rrapi et al.\ survey their emerging roles in reasoning, planning and
interaction in collaborative robotics \cite{ref39}. A proposal can be
reasonable as a description of what to do and still be wrong as a physical
command.

The constraints that decide whether a command is acceptable depend on the
task and on its current state. Removing a screw couples axial motion to
rotation through the thread pitch, so pulling without turning strips the
thread. A snap-fit latch must be deflected before the part can be pulled.
A crank permits axial extraction only in particular rotation phases, and its
transverse limits turn with the workpiece. Pressing a battery beyond its
deformation onset can start an internal short that keeps heating after the
force is removed. These failures differ in kind. Battery deformation
progressing to an internal short and heating is a modeled hazard, whereas
circuit-board fracture, stripped threads and broken latches are failures of
asset integrity. Throughout this paper we measure modeled command
admissibility and task outcomes under declared constraints; these
measurements do not establish physical safety.

Existing approaches address different parts of this problem. Constrained
reinforcement learning optimizes a policy under cost constraints
\cite{ref16,ref17}; Lagrangian methods adapt penalties during learning
\cite{ref18,ref19,ref20}; safety layers correct actions using a constraint
model \cite{ref21}; and recovery policies intervene when estimated risk is
high \cite{ref22}. Control barrier functions can describe nonlinear and
state-dependent safe sets and, under their dynamics and feasibility
assumptions, support forward-invariance guarantees \cite{ref23,ref24}.
Extensions handle higher relative degree \cite{ref40} and time-varying
barriers \cite{ref41} under their stated dynamical conditions.
Table~\ref{tab:submitted-1} distinguishes these formulations. For a governor
placed between a proposer and a robot, two further questions decide what it
can deliver. The first is which constraint representation and task
information the implemented governor actually has. The second is how the
upstream policy adapts to the governor: a restriction added to a policy
trained without it receives different proposals from one present throughout
training.

\begin{table}[!htbp]
\centering\small
\caption{Representative approaches to constrained learning and execution.
Entries describe the cited formulations and their assumptions. Policy learning
and runtime filtering can be combined.}
\label{tab:submitted-1}
\begin{tabular}{@{}L{0.96in}L{2.60in}L{2.75in}@{}}
\toprule
Approach & Constraint and intervention & Scope of the cited formulation\\
\midrule
CPO \cite{ref16} & Trust-region policy updates constrained by expected cumulative costs. & Policy-level cost guarantees depend on CMDP assumptions and update approximations.\\
Lagrangian RL \cite{ref18,ref19,ref20} & Policy and penalty multipliers adapt to cost signals. & Constraint handling occurs during learning; a separate execution filter can be added.\\
Safety layer \cite{ref21} & Learned state-dependent, action-affine constraint model; minimum-change correction. & Accuracy depends on the model; the cited closed form assumes at most one active constraint.\\
Recovery RL \cite{ref22} & A learned risk critic switches execution to a recovery policy. & Depends on risk estimation and the recovery behavior.\\
CBF-QP \cite{ref23,ref24} & Dynamics and barrier derivatives define state-dependent control constraints. & Nonlinear safe sets are allowed; invariance requires the stated regularity, model and feasibility assumptions.\\
TRiX & Task context selects explicit command sets and an execution outcome. & Tests command admissibility and retained authority under declared task models, with frozen or filter-aware proposers.\\
\bottomrule
\end{tabular}
\end{table}

We present TRiX, an execution governor that sits between proposal generation
and actuation (Figure~\ref{fig:submitted-1}). An upstream proposer, such as a
reinforcement-learning policy or a VLM planner, produces a command. The
current task state and a declared task context select the set of commands
that the task model admits. TRiX then passes the command unchanged, projects
it onto that set, rejects it when no meaningful correction exists, or reports
that the selected specification leaves no usable command. Only a passed or
projected command reaches the executor. This separation makes two evaluation
regimes explicit. A governor applied to a frozen policy must handle proposals
generated without adapting to its restrictions. A policy trained with the
governor active can instead reorganize its proposals around the authority
that the governor preserves. Performance in one regime does not establish
performance in the other.

\begin{figure}[!htbp]
\centering
\includegraphics[width=\textwidth]{manuscript/figures/trix_architecture_fig1.pdf}
\caption{Action proposal and execution authority in TRiX. A learned or
generative planner proposes a manipulation command. TRiX evaluates the
proposal against task-conditioned physical admissibility constraints before
the command reaches the execution system.}
\label{fig:submitted-1}
\end{figure}

We evaluate TRiX at three levels, each answering a different question.
DISASM-Bench, a set of six analytical disassembly tasks, supports controlled
comparisons of trained policies, runtime restrictions and constraint
representations under both regimes. Recorded commands from three VLM runs on
shared circuit-board images are replayed with and without governance from
identical states, testing governance of actual model outputs. A physical
B601-RS arm tests whether governance decisions change what reaches the
actuators. The analytical tasks are models with declared parameters; neither
the replay study nor the robot experiments calibrate their force, torque or
damage thresholds.

The experiments show that a governor and the policy it governs must be
evaluated together. On circuit-board extraction, the same componentwise
restriction gives 0.0\% safe completion when applied to an already trained
policy and 100.0\% when present during training. Training does not recover
everything: on crank extraction, a static restriction stays at 0.0\% while
the task-conditioned governor reaches 49.4\%. With the specification matched,
analytical projection and a quadratic program give identical outcomes on
screw extraction, which separates solver choice from command-set
representation. Outside the training reset distribution, snap-fit safe
completion under TRiX falls from 96.5\% to 59.9\%, and latch fractures occur
even though every executed command stays inside its bounds. In replay of
recorded model commands, governance prevents all 31 observed
constraint-violation cases among 257 grounded commands and introduces none,
and on the physical arm it rejects or corrects commands before actuation.
TRiX is therefore best understood as an auditable execution contract whose
value depends on the proposer, the specification and the task context it is
given.

\section{Methods}

\subsection{Execution contract}
\label{sec:authority}
Let $x$ be the observed state, $c$ the declared task context and
$\mathcal A(x,c)$ the applicable admissible command set. A proposer supplies a
command $\tilde a$. TRiX returns \texttt{ALLOW} for a proposal already in that
set, \texttt{PROJECT} for an inadmissible proposal with a feasible correction,
\texttt{REJECT} for a non-projectable symbolic or task-state violation, and
\texttt{INFEASIBLE} when the selected specification supplies no usable
command authority. Rejection and infeasibility do not forward a command for
execution. For a nonempty closed convex set in normalized command coordinates,
the correction is the unique Euclidean projection
\begin{equation}
\Pi_{\mathcal A}(\tilde a)=\arg\min_{a\in\mathcal A(x,c)}
\tfrac12\|a-\tilde a\|_2^2.
\label{eq:projection}
\end{equation}
Normalization is part of the model: squared velocity, angular-velocity and
force errors require declared scales before they can share this metric.

Screw extraction shows why the representation of $\mathcal A$ matters. With
pitch $p$, axial velocity $v_z$ and angular velocity $\omega_z$, ideal motion
obeys $v_z=k\omega_z$, where $k=p/(2\pi)$. Using command scales $S_v$ and
$S_\omega$ gives normalized commands $a_v=v_z/S_v$ and
$a_\omega=\omega_z/S_\omega$, and a coupling coefficient
$\alpha=kS_\omega/S_v$. The projection onto the ideal line is
\begin{equation}
\Pi_0(\tilde a)=
\frac{\alpha\tilde a_v+\tilde a_\omega}{1+\alpha^2}
\begin{bmatrix}\alpha\\1\end{bmatrix}.
\label{eq:helix}
\end{equation}
Independent coordinate bounds do not encode this coupling. The implemented
SCREW task uses a finite-tolerance strip with command and radial-force bounds
(Figure~\ref{fig:submitted-2}); an appropriately specified linear-constraint
quadratic program (QP) represents the same set. This makes solution mechanism
separable from constraint representation.

\begin{figure}[!htbp]
\centering
\includegraphics[width=\textwidth]{manuscript/figures/figure2_representation_gap.pdf}
\caption{Representation of the SCREW command coupling. (a) A two-dimensional
slice of the normalized command box, with radial action fixed at zero.
The frozen environment uses pitch $p=1.25$ mm and the finite tolerance
$|v_z^c-p\omega_z^c/(2\pi)|\leq0.001$ m\,s$^{-1}$; the admissible strip
occupies 2.0\% of this slice. Commands are normalized by 60 rad\,s$^{-1}$
and 0.05 m\,s$^{-1}$, respectively; the displayed coefficient is rounded.
The tightest enclosing componentwise box admits commands outside the strip.
(b) The same illustrative proposal is corrected by componentwise clipping
and by Euclidean projection in normalized action coordinates. The latter
uses the coupled strip and command bounds. The dashed box is a geometric
illustration, not a registered SCREW experimental arm; an appropriately
specified linear-constraint QP also represents the strip.}
\label{fig:submitted-2}
\end{figure}

Four representation properties frame the analysis. First, projection onto an
axis-aligned product of intervals is componentwise clipping. Second, a
phase-independent set valid in two phases must lie in their intersection,
which may discard phase-specific authority. Third, the intersection over all
rotations of a rectangle with half-widths $a\geq b>0$ is the radius-$b$ disk;
a fixed square $[-h,h]^2$ valid for every orientation therefore requires
$h\leq b/\sqrt2$. Fourth, a Gaussian-tightened scalar interval
$|y|\leq L-z\sigma$ retains nonzero authority exactly when $z\sigma<L$.
The assumptions and proofs are in \suppsection{S1}, \suppsection{S2} and
\suppsection{S3}. These are statements about representation and feasibility.
They do not imply closed-loop stability or forward invariance for an
unmodeled plant. How a trained policy adapts to a restriction is examined
empirically rather than stated as a proposition.

\subsection{Disassembly tasks and outcome definitions}
\label{sec:tasks}
DISASM-Bench supplies six analytical environments: SCREW (screw extraction),
PCB (circuit-board extraction), SNAP (snap-fit release), CRANK (crank
extraction), BATTERY (battery peel) and PRY (prying a bonded part open with a lever). Each
maps three normalized action coordinates to a task-specific physical
interface (Table~\ref{tab:submitted-2}). The commanded quantities are axial
and angular velocity $v_z^c,\omega_z^c$ and radial force $F_r$ for SCREW;
tilt torques $\tau_x,\tau_y$ and lift force $F_z$ for PCB; latch-deflection,
pull and lateral forces $F_{\rm def},F_{\rm pull},F_{\rm lat}$ for SNAP;
forces $F_x,F_y,F_z$ for CRANK; peel force $F_{\rm peel}$ and lateral forces
$F_{\ell1},F_{\ell2}$ for BATTERY; and tool torque $\tau$, insertion force
$F_{\rm ins}$ and lateral force $F_{\rm lat}$ for PRY, whose torque limit
depends on insertion depth $d$.

\begin{table}[!htbp]
\centering\small
\caption{Task interfaces and nominal modeled boundaries. Limits are environment
specifications, before chance tightening. The three normalized actions map to
the listed physical commands; scales, observations, damage accumulation and
reset distributions are in \suppsection{S4}.}
\label{tab:submitted-2}
\begin{tabular}{@{}L{0.74in}L{1.30in}L{2.41in}L{1.72in}@{}}
\toprule
Task & Command & Constraint or context & Damage/progress mechanism\\
\midrule
SCREW v3 & $v_z^c,\omega_z^c,F_r$ & Pitch $p=1.25$~mm; $|v_z^c-p\omega_z^c/(2\pi)|\leq10^{-3}$ m\,s$^{-1}$; $|F_r|\leq20$~N. & Thread engagement, 35-N load capacity and relative slip.\\
PCB v2 & $\tau_x,\tau_y,F_z$ & $\|\symbfit\tau_{xy}\|\leq0.100$~N\,m; $|F_z|\leq40$~N. & Bending damage above 0.05~rad; fracture above 0.075~rad or damage~1.\\
SNAP v2 & $F_{\rm def},F_{\rm pull},F_{\rm lat}$ & Engaged pull $\leq5$~N; release at 2-mm deflection or retained release; lateral bound 15~N. & Premature-pull damage; breakage beyond 6-mm deflection or damage~1.\\
CRANK v2 & $F_x,F_y,F_z$ & Task-frame torque $\leq1.5$~N\,m and radial force $\leq20$~N; axial bound changes with rotation phase. & Angle-dependent rotation/extraction and accumulated damage.\\
BATTERY v2 & $F_{\rm peel},F_{\ell1},F_{\ell2}$ & Flagged peel bound 20~N; lateral norm 15~N; temperature 60$^\circ$C. & Deformation begins above 15~N; accumulated deformation can initiate a short and delayed heating.\\
PRY v2 & $\tau,F_{\rm ins},F_{\rm lat}$ & $|\tau|\leq1.5[1-0.5\min(1,|d|/0.05)]$ N\,m; insertion 50~N; lateral 20~N. & Depth-dependent crack accumulation, bond release and gap growth.\\
\bottomrule
\end{tabular}
\end{table}

The versioned registry binds SCREW v3 and the five v2 tasks to their
analytical environments and named filters. All observations contain ten
float32 entries; their meanings are task-specific, and unused entries are
zero padding. For example, PCB contains lift and lift velocity, two tilt
angles and their rates, accumulated damage and a fracture flag. Observations
are not one shared position/velocity/wrench vector. Actions are normalized to
$[-1,1]^3$ and mapped through the physical scales in \suppsection{S4}.
Each outer action advances $1/240$~s of simulated time. SCREW uses 40
internal integration substeps; PCB, SNAP, CRANK and BATTERY use ten; PRY
updates directly. Progress and termination are task-specific.

Safe completion combines successful task termination with the absence of the
declared outcome-affecting violation. The outcome contract also distinguishes
completion with outcome-affecting violations, destructive completion and the
remaining unsuccessful outcomes; diagnostic events are recorded separately.
Reset parameters, distributions, ranges and training/evaluation use are
tabulated in \suppref{tab:submitted-9} and detailed in \suppsection{S4.6}.
Risk allocations, noise and reset variation come from the frozen source.

Figure~\ref{fig:disasm-workspaces} visualizes recorded states of the six
tasks in PyBullet \cite{ref33}, with a per-panel state manifest. The
analytical environments generate the reported dynamics and outcomes; the
renderer is not a second validation of the analytical equations or a
simulated articulated-arm controller, and the illustrated gripper does not
imply that articulated-arm or contact dynamics generated the policy results.
An earlier PyBullet helical-constraint prototype imposed the same screw pitch
through a gear constraint. It therefore could not independently validate that
pitch; we withdrew the associated validation argument.

\begin{figure}[!htbp]
\centering
\includegraphics[width=\textwidth]{manuscript/figures/disasm_bench_workspaces.pdf}
\caption{Six recorded disassembly workspaces. (a) SCREW: coupled threaded
extraction. (b) PCB: lift and two-axis board tilt, with $\|\symbfit\theta\|=
0.00748$ rad in this frame. (c) SNAP: latch deflection before extraction.
(d) CRANK: rotation and axial removal. (e) BATTERY: peel through an upright
extraction tab. (f) PRY: opening one end with the opposite edge supported;
the orange element is the tool handle. Poses come from recorded environment
states and the per-panel manifest. RGB axes mark world-aligned display
references; grid spacing is 20 mm. Purple arrows are schematic command/load
directions, not measured force or torque. PyBullet supplies visualization;
arm/contact dynamics are not simulated. Panels are individual snapshots.}
\label{fig:disasm-workspaces}
\end{figure}

\subsection{Policies and evaluation regimes}
\label{sec:training-methods}
We train Proximal Policy Optimization (PPO) \cite{ref31} and Soft
Actor-Critic (SAC) \cite{ref32} policies with Stable-Baselines3 2.9.0,
Gymnasium 1.3.0 and two-layer $[256,256]$ networks on CPU. Ten independent
training seeds, 0--9, define each policy comparison. Stage~1 trains without
the evaluated governor and then applies each runtime arm in
Table~\ref{tab:submitted-3} to the same selected frozen policy. Stage~2 trains
a separate policy with its restriction inside the environment transition;
replay and rollout buffers retain the nominal policy action. The adapter
distinguishes termination from time-limit truncation and restores NumPy
random state after interleaved evaluation.

At each predeclared task budget, checkpoints saved every 1,000 steps are
screened on 30 fixed validation episodes. Three-checkpoint windows are ranked
by mean safe completion, and the centers of the three best windows receive
70 additional validation episodes. Selection is lexicographic: maximize safe
completion, minimize destructive completion, maximize normalized progress,
minimize episode-level violation incidence, then choose the earliest
checkpoint. Each selected policy is tested on 100 held-out episodes, disjoint
from validation; pilot seeds are excluded. Recorded checkpoint lists,
software fingerprints, task budgets and episode-seed exceptions are retained
in \suppsection{S5} and \suppsection{S6}.

\subsection{Runtime arms and specification controls}
\label{sec:controls-methods}
The runtime registry distinguishes componentwise restrictions, analytical
local-gradient references, matched QP and task-specific TRiX operators
(Table~\ref{tab:submitted-3}). A learned constraint-representation experiment
is evaluated separately from these policy results. No trained PPO-Lagrangian
or complete trained SafeLayer policy is included in the quantitative matrix.
Literature comparisons in Table~\ref{tab:submitted-1} retain their original
assumptions and are not relabeled as experimental baselines.

\begin{table}[!htbp]
\centering\small
\caption{Trained proposers, runtime arms and representation controls.
Stage~1 freezes the policy before applying an arm; Stage~2 trains with the arm
inside the transition. A method's name does not imply a shared specification
across tasks.}
\label{tab:submitted-3}
\begin{tabular}{@{}L{1.15in}L{2.66in}L{2.50in}@{}}
\toprule
Entry & Implemented operation & Experimental role\\
\midrule
PPO / SAC \cite{ref31,ref32} & Learn nominal proposals from task observations and reward. & Only trained task-policy algorithms reported.\\
\texttt{none} & Pass the proposal through. & Unfiltered frozen-policy reference.\\
\texttt{box\_clip} & Clip coordinates independently. & PCB and BATTERY; preventive BATTERY variant uses a matched 15-N target.\\
\texttt{static\_clip} & Use the task's context-limited componentwise bounds. & SNAP ignores latch phase; PRY uses zero-insertion bounds; CRANK retains axial phase but fixed transverse axes.\\
\texttt{oracle\_tangent} & Iterate analytical local-gradient halfspace corrections and bound clipping. & Task-specific reference in Stage~1; no learned constraint model.\\
\texttt{qp\_matched} & Active-set QP over the declared SCREW command polytope. & Matched solution-mechanism control; no additional state-space barrier certificate.\\
\texttt{trix} & Task-specific analytical projection or bounded correction. & All six tasks; BATTERY Stage~2 uses \texttt{trix\_preventive}.\\
Learned affine / nonlinear \cite{ref21} & Fit constraint models and correct using their gradients. & Separate representation studies; not trained SafeLayer policy baselines.\\
\bottomrule
\end{tabular}
\end{table}

The SCREW matched QP and analytical solver minimize the same normalized
Euclidean error over the same finite helical slab, action bounds and radial
limit. This controls command-set specification; it does not construct a
state-space barrier certificate. BATTERY preventive arms share the 15-N
threshold and margin construction. PCB and SNAP comparisons retain their
implemented margin differences, and CRANK combines frame and angle
conditioning. Historical BATTERY prefix pooling and the PCB oracle
episode-seed correction are documented alongside preserved originals. The
full policy tables use exact-arm counts; no checkpoint was reselected during
that audit.

\subsection{Representation, sensitivity and prospective tests}
\label{sec:model-validation}
The PCB decomposition uses the archived seed-0 run with 20,000 fitting
samples and 2,000 held-out proposals, plus the 500--20,000-sample
affine-model sweep. The procedural generator defines six three-dimensional
set families, with 120 training, 40 validation and 200 held-out geometries
from disjoint seed streams. Learned constraint models are fitted separately
on each evaluated instance using 3,000 labeled state-action samples;
nonlinear fitting uses 250 gradient steps. Each arm receives 400 sampled
proposals per instance. Required-action retention means displacement below
$0.02(\|u_{\rm req}\|+10^{-9})$. The context-free box's required-action
exclusion defines the 124/76 partition. Approximate arms share implemented
inscribed-box bounds, which can contribute both feasibility and conservatism.

The sensitivity study varies 21 constants across five multipliers from 0.5
to 2.0, with eight scripted episodes per arm and case. Its ordering criterion
first compares destructive completions and, when both are zero, safe
completion. The nine failed cases are diagnosed after the sweep. The separate
BAYONET experiment preregisters a 150,000-step budget, SAC as the primary and
PPO as a replication algorithm, and a 50\% competence gate; neither exceeds
10\% mean safe completion. The gate and task are not changed after this
inconclusive result. \suppsection{S7} and \suppsection{S10} give
implementation checks, designs and full outcomes.

\subsection{Shifted reset evaluation}
\label{sec:ood-methods}
Held-out tests in the original evaluation use the training reset ranges.
A prospectively declared evaluation adds 42,000 episodes across the five
tasks with stochastic physical resets. It widens existing reset half-widths
by $1.5\times$ and $2\times$ and samples a $2\times$ outside-support shell,
with selected policies and governors fixed. PRY is not applicable because its
physical reset is deterministic. \suppsection{S14} and \suppref{tab:ood-ranges}
give the committed protocol, explicit reset interface and complete ranges.
The SNAP command and damage attribution reported in Results was conducted
after this evaluation and is labelled post hoc.

\subsection{Recorded vision-language-model commands}
\label{sec:vlm-method}
The frozen bank has 28 base configurations: 21 have normal, rotated,
occluded and misleading views, and seven have one partial view each,
producing 91 images. Each provider receives the same image/prompt pairs and
a 2,000-token output limit. Recorded aliases are
\texttt{claude-sonnet-4-5}, \texttt{gpt-4o} and
\texttt{gemini-flash-latest}; they are not pinned backend snapshots
\cite{ref36,ref37,ref38}. All raw responses and parsed commands are retained.
Backend IDs, per-call dates and retry attribution are unavailable. Four
retried Gemini timeouts are mentioned in historical notes, but their image
IDs and successful variants cannot be recovered; the client also permits
fallback outside those events.

The images use a visibility-scaled tilt cue. Main-stratum display tilt is
approximately $20.669$ times replayed mechanical curvature; partial views use
$3.6$ times. Saved near-boundary mechanical tilts of 0.0358--0.042525~rad
appear as roughly $42.4^\circ$--$50.4^\circ$, beyond the model's $0.075$-rad
($4.30^\circ$) fracture threshold if interpreted physically. Thus the
displayed pose need not be a realizable intact-board state. Positive scaling
preserves ordering within the 84 main-stratum images, but not physical
distance to a boundary; all margin comparisons exclude the seven partial
views. Rendering reproducibility does not establish fidelity to mechanical
tilt. The stratum counts describe responses to this encoding. Without an
unscaled image control, effects of visual scaling cannot be separated from
differences between model runs.

The fixed parser maps outputs to \texttt{LIFT}, \texttt{TILT\_X},
\texttt{TILT\_Y}, \texttt{PRESS}, \texttt{HOLD} or \texttt{ABORT}.
Refusals and ungroundable outputs are distinct from grounded admissible and
inadmissible commands. Set membership defines inadmissibility; no independent
semantic-correctness annotation is available. Each grounded command is held
for 40 steps in unfiltered and governed replay from the same restored state
and seed. A violation case contains at least one recorded constraint
violation; prevention requires violations only in the unfiltered arm.
\suppsection{S11} provides the display formulas, trial ledger, stratum counts
and conditional Gemini sensitivity bound.

\subsection{Physical platform}
\label{sec:hardware-methods}
The Seeed reBot B601-RS uses RobStride actuators, a 48-V supply and
1-Mbit/s CAN. Joints~1--5 use 0.1-second MIT-style position/velocity
commands; Joint~6 uses an approximately 200-Hz velocity controller. The final
grasp implementation is identified by commit \texttt{d296e6e} and tag
\path{trix-hardware-soft-grasp-success-20260828}. TRiX separates symbolic
proposals, decision receipts and approved executor targets, with Joint~1
protected by a hard invariant. The calibrated sequence uses J3 gravity
feed-forward and a disclosed Joint~2 positioning offset around engagement.
In all six grasp-and-return trials, the operator manually positions the
object between the open jaws at the calibrated grasp-ready pose before
governed closure, lift and return. The placement protocol for the three
repetitions without video is author-confirmed; their logs record robot
execution. Console logs, passive CAN captures, source/configuration
snapshots, video mappings and checksums support the intervention experiment,
three video-backed trials and three logged repetitions. There is no
independent force/torque measurement or autonomous perception-to-grasp
evaluation (\suppsection{S12}).

\subsection{Runtime measurement}
\label{sec:timing-methods}
The latency benchmark times the complete $(o,a)\mapsto(a',\mathrm{info})$
call, including allocations, per-call constraint construction and
diagnostics. It uses 2,048 saved inputs per task, five fresh processes and
256 warm-up calls per implementation/process. Identical per-task input
sequences are shared across arms, with task/arm order randomized in 64-input
blocks. Durations are untrimmed integer nanoseconds; cyclic garbage
collection is disabled, and timer overhead is recorded without subtraction.
Each process uses CPU~0 and one native numerical thread. OSQP reuses its
problem and warm start; one-time construction is reported separately.
Policy/VLM inference, input generation, simulation and robot communication
are outside the interval. \suppsection{S8} records host/software versions,
tolerances, raw arrays, source hashes and the complete workload design.

\subsection{Statistical analysis}
\label{sec:statistics}
The training seed, not the evaluation episode, is the unit for policy
comparisons. We report ten-seed means and all plotted seed values. Selected
contrasts use paired differences by seed identifier and exact percentile
bootstrap distributions over those ten differences, computed by convolution.
Intervals are descriptive, not multiplicity-adjusted tests. Same-policy
Stage~1 contrasts are distinguished from separately trained Stage~2 and
cross-regime pairs. The preventive BATTERY and PCB cross-regime contrasts
also have recorded episode-seed mismatches, so pairing is only at
training-seed level (\suppref{tab:submitted-11}). The shifted reset study
uses the same estimator and training-seed pairing. Its changes from
historical in-range estimates include episode-sampling variation because the
schedules differ; within/outside-support strata are descriptive
(\suppsection{S14}). Procedural instances are the units in the geometry
study. VLM counts are descriptive over paired images and commands; views
clustered within 28 configurations are not independent physical replicates,
and no inferential cross-model test is reported.

\subsection{Records and provenance}
\label{sec:provenance}
Every experimental record is linked to the applicable environment, policy,
constraint, margin and outcome specifications. The curated source map links
all main and supplementary figures and tables to data or executable source.
Correction notes retain affected historical files. Independent physical
calibration and unresolved backend identities remain provenance limits.
Figure~\ref{fig:seed-distributions} was regenerated from the retained
seed-level data using a supplied replacement plotting script; the original
plotting script remains unavailable.

\section{Results}

\subsection{Training with the governor changes which restriction suffices}
\label{sec:policy-results}
The clearest regime effect occurs on PCB extraction. When the same
componentwise \texttt{box\_clip} restriction is applied after SAC training,
safe completion is 0.0\%; with that restriction active during training, it is
100.0\% (Table~\ref{tab:policy-summary}). Every one of the ten training seeds
shows this change (Figure~\ref{fig:seed-distributions}a). The same frozen
policy reaches 100.0\% under TRiX, so the post-hoc failure reflects the
restriction rather than an incompetent policy. A policy trained under the box
learns to use the authority the box retains. This result establishes the role
of adaptation while holding the implemented restriction fixed.

\begin{table}[!htbp]
\centering\small
\caption{Key policy comparisons within the training reset ranges. Values are
mean safe completion (\%) across 10 training seeds, with 100 held-out episodes
per selected policy/arm.
Stage~1 rows compare the same frozen policy; Stage~2 rows compare separately
trained policies. Full arm tables and seed-paired intervals are in
\suppsection{S6} and \suppsection{S9}.}
\label{tab:policy-summary}
\begin{tabular}{@{}llllrr@{}}
\toprule
Task & Policy & Regime & Compared restrictions & First & Second\\
\midrule
SCREW & SAC & Stage 1 & Matched QP / TRiX & 88.7 & 88.7\\
SCREW & PPO & Stage 1 & Matched QP / TRiX & 60.0 & 60.0\\
PCB & SAC & Stage 1 & Box / TRiX & 0.0 & 100.0\\
PCB & SAC & Stage 2 & Box / TRiX & 100.0 & 99.8\\
SNAP & SAC & Stage 2 & Static / TRiX & 5.2 & 96.5\\
CRANK & SAC & Stage 2 & Static / TRiX & 0.0 & 49.4\\
PRY & SAC & Stage 2 & Static / TRiX & 88.5 & 89.1\\
BATTERY & SAC & Stage 2 & Preventive box / TRiX & 100.0 & 100.0\\
\bottomrule
\end{tabular}
\end{table}

\begin{figure}[!htbp]
\centering
\includegraphics[width=\textwidth]{manuscript/figures/figure4_seed_distributions.pdf}
\caption{Seed-level safe completion within the training reset ranges in the
frozen production experiments.
Each point is one independent training seed (\(n=10\)). Boxes span the
interquartile range, central lines show medians, whiskers extend to the most
extreme observations within 1.5 interquartile ranges of the box, and triangles
denote arithmetic means.
(a) PCB compares the same \texttt{box\_clip} restriction under post-hoc
Stage~1 governance and Stage~2 filter-aware training, yielding mean safe
completion of 0.0\% and 100.0\%, respectively.
(b--d) Stage~2 filter-aware SAC comparisons for SNAP, CRANK, and PRY.
SNAP yields 5.2\% under the static restriction and 96.5\% under the implemented
task-conditioned governor. CRANK yields 0.0\% under the static restriction and
49.4\% under the task-conditioned governor, with the latter exhibiting a
bimodal seed distribution. PRY yields similar distributions under the simpler
and task-conditioned restrictions, with means of 88.5\% and 89.1\%.}
\label{fig:seed-distributions}
\end{figure}

\subsection{Matched controls separate solver from representation}
\label{sec:matched-results}
SCREW holds the specification fixed and varies only the solution mechanism.
Under frozen-policy evaluation, \texttt{oracle\_tangent}, \texttt{qp\_matched}
and TRiX each achieve 88.7\% safe completion for SAC and 60.0\% for PPO,
against 0.0\% without a governor. The matched constrained controls reach
100\% for both algorithms with filter-aware training. Equal outcomes are
expected when the same admissible command set is supplied to each method;
on this task, differences between methods arise from what each is given,
not from how it computes the correction. Implementation cost is reported
separately below.

\subsection{Missing task context limits what training can recover}
\label{sec:context-results}
Adaptation does not recover all lost authority. CRANK's static restriction
remains at 0.0\% after filter-aware SAC training, while the implemented
task-conditioned governor reaches 49.4\%. That mean conceals seeds near
0\% or 100\% (Figure~\ref{fig:seed-distributions}c). The comparison changes
both angle-conditioned axial bounds and transverse task-frame representation,
so it does not isolate either factor.

SNAP separates as strongly. Within the training reset ranges, SNAP reaches
5.2\% under static and 96.5\% under TRiX
(Figure~\ref{fig:seed-distributions}b). This compares separately trained
policy-and-governor systems and does not isolate the governor's causal
contribution; the implementations also differ in margin and latch-phase
conditioning. The shifted reset evaluation below reduces these rates to
3.5\% and 59.9\%, respectively.

\subsection{Simpler restrictions suffice when context and specification match}
\label{sec:sufficient-results}
On PRY, filter-aware SAC gives 88.5\% and 89.1\% under static and
task-conditioned restrictions, with closely aligned seed distributions and
the same zero-performing seed (Figure~\ref{fig:seed-distributions}d). PCB
reaches 100.0\% with the box and 99.8\% with TRiX in this regime. The
production PCB comparison also changes the torque margin: the box uses a
nominal 0.100~N\,m per-axis bound, while TRiX uses a robust coupled radius of
approximately 0.087087~N\,m. Consequently, its between-arm comparison does
not isolate box-versus-disk geometry.

BATTERY isolates specification. Historical arms compared the 20-N flagged
peel boundary with a 15-N preventive deformation threshold. When both
componentwise and TRiX arms use the same preventive target and margin
construction, both reach 100\% safe completion. Historical aggregation also
pooled similarly named BATTERY arms; the exact-arm audit in
\suppsection{S9.4} corrects Stage~1 SAC/TRiX safe completion to 20.0\%, with
original records retained. \suppsection{S9} reports the full policy tables
and these controls; \suppref{tab:submitted-11} supplies training-seed
intervals and pairing exceptions. Nominal-policy failure on some tasks leaves
frozen-policy arms at zero, which limits the ranking those cells can support.

\subsection{Learned constraint models can lose required authority}
\label{sec:representation-results}
The PCB representation study distinguishes learned constraint models from
trained task policies. On 2,000 held-out proposals, clipping produces 2.9\%
feasible commands, a learned action-affine model 0.0\%, a learned nonlinear
model 78.0\% after one correction and 92.3\% after iteration, and the
analytical oracle and exact projector 100\%. The affine result remains 0.0\%
when fitting samples increase from 500 to 20,000. This single-seed study
tests the implemented representations; it does not establish a general
failure of learning or an ordering of complete safety algorithms.

A separate procedural study evaluates 200 admissible-set instances from box,
ellipsoid, polytope, state-scaled, rotating and phase-gated families. Each
has a required action sampled inside its exact set. The implemented
context-free box excludes that action in 124 instances and retains it in 76.
This partition is a constructed calibration: the box's exclusion is the
grouping criterion, so its behavior cannot independently validate that
criterion. The alternatives reveal information beyond the partition
definition (\suppref{tab:procedural-geometry}). Exact projection retains the
required action in 100\% of both groups; the context-conditioned oracle
retains it in 36.3\% and 100\%, respectively. The learned affine model
retains it in 0\% of both groups, including the group where a box succeeds.
All arms return feasible commands on 400 sampled proposals per instance,
partly because approximate arms also use the implemented inscribed-box
bounds. Feasibility can therefore coexist with loss of task-required
authority. Models are fitted separately on samples from each evaluated
geometry; this is not zero-shot transfer of one learned constraint model
(\suppsection{S10}).

\subsection{Admissible commands do not guarantee safe outcomes outside the training range}
\label{sec:ood-results}
In the primary outside-training-support shell of the prospectively declared
evaluation, SNAP safe completion is 59.9\% under TRiX versus 3.5\% under
static. These rates fall from the in-range 96.5\% and 5.2\%. As in range,
this compares separately trained policy-and-governor systems and does not
isolate the governor's causal contribution. Latch fractures, including
timeouts, occur in 401 versus 468 of 1,000 episodes per arm. Destructive
completions number 369 versus 886; this category combines completed episodes
with fracture or accumulated damage and excludes timeouts
(\suppref{tab:snap-ood-attribution}). The advantage in safe completion is
therefore much larger than the reduction in fractures.

A post-hoc replay audit finds every TRiX command inside the frozen executable
bounds. Nevertheless, all 401 TRiX fractures begin through overdeflection in
the softer-latch stratum, before pull damage; all 480 TRiX episodes in the
harder-latch stratum complete safely. Per-step command membership therefore
did not secure the subsequent dynamic state. These observed counts do not
establish safety throughout either stiffness interval. \suppsection{S14.4}
reports the full outcome partition, traces and the distinction between
executed margins and the later, tighter registry.

The other tasks change less. CRANK reaches 47.6\% under TRiX versus 0.0\%
under the static arm, with 22 destructive completions per 1,000 episodes
under TRiX. Matched SCREW arms remain equal (87.9\% SAC; 59.9\% PPO). PCB
yields 0.0\%/99.6\% in Stage~1 and 100.0\%/99.8\% in Stage~2 for box/TRiX;
preventive BATTERY gives 100.0\%/100.0\%. \suppsection{S14} and
\suppref{tab:ood-safe} report all three declared conditions, changes and
intervals; \suppref{tab:ood-outcomes} reports every outcome category. The
original in-range results remain unchanged.

\subsection{Governance intercepts violations in recorded model commands}
\label{sec:vlm-results}
Three recorded provider runs receive the same 91 PCB images and prompt,
yielding 273 decisions. The labels Claude Sonnet, GPT-4o and Gemini identify
the configured aliases and recorded runs; resolved backend versions are
unknown. Table~\ref{tab:vlm-results} gives the complete decision partition:
16 refusals, 219 admissible commands and 38 inadmissible commands, with no
ungroundable outputs. Both admissible and inadmissible commands are grounded
and eligible for paired replay, giving 257 command pairs.

\begin{table}[!htbp]
\centering
\caption{VLM command study on 91 shared images per run, using the configured
model labels defined in Methods. Refused, admissible and
inadmissible counts partition all decisions; ungroundable outputs are zero.
Grounded equals admissible plus inadmissible. A violation case is a replay
with at least one constraint violation; a prevented case has violations only
in the unfiltered arm. All governed violation counts are zero. The last two
columns count command cases, not violating timesteps or damaged workpieces.}
\label{tab:vlm-results}
\begingroup
\footnotesize
\setlength{\tabcolsep}{2pt}
\begin{tabular}{lrrrrrr}
\toprule
Model run & Refused & Admissible & Inadmissible & Grounded &
\shortstack{Unfiltered\\violation cases} &
\shortstack{Prevented\\violation cases}\\
\midrule
Claude Sonnet & 0  & 88 & 3  & 91 & 0  & 0\\
GPT-4o        & 16 & 66 & 9  & 75 & 9  & 9\\
Gemini        & 0  & 65 & 26 & 91 & 22 & 22\\
\midrule
Total         & 16 & 219 & 38 & 257 & 31 & 31\\
\bottomrule
\end{tabular}
\endgroup
\end{table}

TRiX prevents all 31 observed constraint-violation cases and introduces none.
A prevented case has at least one violation during the unfiltered replay and
none under governance from the same restored state and seed; merely changing
a command receives no prevention credit. The three inadmissible Claude
commands produce no unfiltered violation cases; GPT-4o contributes nine
prevented cases and Gemini 22. For example, GPT-4o's \texttt{s06\_occluded}
response requests \texttt{LIFT} at magnitude 0.8, grounding to $(0,0,0.8)$:
violations occur on 40/40 unfiltered steps and 0/40 governed steps. Both arms
in every replay finish the 40-step horizon without task completion or
mechanical failure. The endpoint is prevented constraint violations, not
fractures or successful extraction. The commands were elicited from
visibility-scaled images whose displayed tilts can exceed the model's
intact-board range, so this result measures governance of those recorded
commands under mechanical replay.

Thirty-five images elicit an inadmissible proposal in at least one run; none
does so in all three. Claude's three inadmissible cases and GPT-4o's nine do
not overlap, even when grouped by base configuration. Holding those responses
fixed, no Gemini failure set can create a three-way overlap. The study
measures one proposal per image, with repeated views of 28 configurations,
and makes no claim about an autonomous agent loop or hallucination detection.

\subsection{Governance changes what reaches a physical robot}
\label{sec:hardware-results}
On the B601-RS, the intervention experiment rejects an invalid symbolic
transition and a proposed $+50^\circ$ motion of protected Joint~1, producing
no physical target. An excessive Joint~2 proposal of $+50^\circ$ is projected
to $+18^\circ$; only the corrected target is transmitted and executed. The
allowed return finishes with Joint~1 drift of $+0.007^\circ$.

The frozen final implementation then completes three video-backed
grasp-and-return trials with the object retained
(Figure~\ref{fig:b601-hardware}). Each contains ten \texttt{ALLOW} decisions,
for 30 governed decisions in total; these trials are not additional
projection or rejection interventions. The maximum absolute Joint~1 return
drift is $0.027^\circ$ and the largest start-to-final residual over
Joints~1--6 is $0.122^\circ$. Three further logged repetitions complete
normally but are not counted as video-backed object-retention trials. The
grasp pose is calibrated. In all six trials, the operator manually positions
the object between the open jaws; the governed sequence then closes the
gripper, lifts and returns. No independent wrench sensor is used. This
demonstrates physical execution of the command boundary, while physical
calibration of the six simulated task models remains open
(\suppsection{S12}).

\begin{figure}[!htbp]
\centering
\includegraphics[width=\textwidth]{manuscript/figures/b601_hardware_sequence.jpg}
\caption{Physical B601-RS execution under the frozen V4 implementation.
(a) Governed approach with the gripper open. (b) The operator places the
object between the open jaws at the calibrated grasp-ready pose immediately
before engagement. (c) Object retained during automatic return. (d) Returned
configuration with the object still held. The grasp pose is calibrated rather
than perception-generated; the experiment evaluates physical
execution-governance integration rather than autonomous grasp synthesis.}
\label{fig:b601-hardware}
\end{figure}

\subsection{Runtime cost depends on the implementation}
\label{sec:runtime-benchmark}
A separate benchmark measures 225,280 complete-filter calls across 22
existing implementations, six tasks and five fresh processes. On the recorded
Intel Core Ultra 9 275HX host, matched SCREW median/95th-percentile latency is
1.255/1.556~$\mu$s for production scalar TRiX, 12.578/15.534~$\mu$s for
active-set QP and 14.924/19.895~$\mu$s for reused OSQP. All matched arms pass
the declared numerical agreement checks on the saved bank; the maximum OSQP
action-coordinate difference is $1.013\times10^{-6}$. The alternative NumPy
analytical implementation takes 18.163~$\mu$s by median, exceeding both QP
medians despite representing the same geometry. The comparison therefore
concerns these implementations and this workload. Across the six registered
TRiX arms, medians range from 1.123 to 11.770~$\mu$s.
\suppref{tab:runtime-latency} reports every implementation, process variation
and the measurement scope. All outliers are retained; these measurements are
not a worst-case or end-to-end robot latency guarantee.

\section{Discussion}
\label{sec:discussion}
The governing restriction and the proposal policy form a joint system. The
PCB cross-regime comparison shows why that matters: a restriction that fails
after training can support successful behavior when included during
training. CRANK shows the complementary limit: adaptation cannot reliably
restore action authority or task context that the supplied restriction
omits. Post-hoc governance and filter-aware learning should therefore be
reported separately, rather than treating one as evidence for the other.

The controls also sharpen how representation should be judged. Matching the
command set makes analytical and QP controls coincide on SCREW, and matching
the preventive specification removes the BATTERY performance difference. The
procedural study distinguishes feasible output from retained required
authority, and the learned nonlinear PCB correction shows that approximation
quality can improve substantially with an expressive representation and
iteration. The contribution is an auditable execution contract and an
experimental account of the conditions under which different implemented
restrictions suffice. It does not establish a universal advantage of symbolic
projection over learned or classical safety methods. The timing comparison
likewise shows an advantage for the scalar implementation on the saved
workload, not an inherent speed advantage of analytical projection over QP:
the NumPy analytical implementation is slower than both matched QP controls.

Several results mark the limits of reactive governance. Outside training
support, every TRiX command on SNAP stayed inside its executable bounds, yet
fractures fell only from 468 to 401 per 1,000 episodes: per-step command
membership does not secure the subsequent dynamic state. In BATTERY, an
internal short can continue heating after external force is removed; a
current-command filter cannot remove accumulated internal state. In SNAP,
suppressing premature pull does not create the latch-deflection action
needed for progress. A correct projector can also enforce the wrong boundary:
the 20-N flagged BATTERY limit allows deformation that begins at 15~N.
Finally, excessive probabilistic tightening can remove all usable authority.
For the reported CRANK axial family, $L/\sigma=5/1.235\simeq4.05$, whereas
the declared multiplier is approximately 4.488. This is reported as a
specification-feasibility failure, with the relevant code path and
assumptions in \suppsection{S13}.

The sensitivity study reinforces that these boundaries are empirical. Its
prescribed task-specific ordering holds in 96 of 105 perturbation cases; the
nine exceptions comprise six PCB boundary/model cases, one CRANK feasibility
failure and two BATTERY completion trade-offs with zero destructive
completions in both arms (\suppsection{S7}). A prospective BAYONET test is
inconclusive because neither policy algorithm reaches the predeclared
competence gate. These results remain part of the evidence rather than being
removed to produce a cleaner narrative.

The VLM and hardware experiments extend the contract to generative proposals
and physical actuation, each within its own scope. The VLM result
demonstrates prevention of recorded constraint violations under paired
command replay, on visibility-scaled images of one task. The robot trials
show that rejection and projection alter what reaches the actuators and that
the same boundary supports repeated manipulation, with a calibrated grasp,
manual object placement and a separate joint-space contract. Neither
calibrates the benchmark's force, torque or damage thresholds. Such
validation would require independent physical measurements and trials
covering contact uncertainty, object variation and model mismatch. Future
evaluation should combine these components with independently calibrated
constraints, state estimation and sequence planning, while measuring which
interventions preserve useful action authority. The supported conclusion is
that explicit execution governance is most informative when evaluated with
its proposal distribution, available context and physical specification made
visible.

\section*{Author Contributions}
Stavan Dholakia: Conceptualization, Methodology, Software, Validation, Writing
(original draft). Shivani Shukla: Formal analysis, Validation, Writing
(review \& editing). Abhishek Singh: Formal analysis, Visualization.
Aditya Gazta: Validation, Visualization, Writing (review \& editing).

\section*{Data Availability}
The source package includes compact policy summaries, seed-level figure
data, VLM responses and trial ledger, procedural and sensitivity results,
workspace state manifests and raw timing records. Checkpoints, episode shards,
video and CAN evidence are checksum-indexed in the research/hardware archive,
preserving original records affected by audit corrections. The package is in
unpublished Zenodo draft 23031294; the editor correspondence supplies its
confidential read-only reviewer link.

\section*{Code Availability}
The Zenodo reviewer link provides \texttt{TRIX\_manuscript\_source.zip}:
manuscript and Supplementary sources, environment/filter implementations,
training and analysis entrypoints, and evidence mappings. Public release of
\url{https://github.com/stavanio/trix-disasm-suite} is planned at resubmission.
Hardware source is identified by \url{https://github.com/stavanio/bhujab601}
and the Methods commit/tag. Supplementary Information records reproduction
protocols and the computational environment.

\section*{Acknowledgements}
None.

\section*{Funding}
This material is based upon work supported by the National Science Foundation
under Grant No. [number] (to [author initials]).
Any opinions, findings, and conclusions or recommendations expressed in this
material are those of the authors and do not necessarily reflect the views of
the National Science Foundation.

\section*{Competing Interests}
S.D. is employed by Microsoft, A.S. by Google and A.G. by Amazon. These
employers provided no funding and had no role in study design, execution,
analysis, writing or the decision to publish.

\clearpage
\begin{thebibliography}{99}

\bibitem{ref1} Forti, V., Baldé, C. P., Kuehr, R. \& Bel, G. The Global E-waste Monitor 2020: Quantities, flows and the circular economy potential. UNU/UNITAR, ITU \& ISWA (2020).

\bibitem{ref2} Harper, G. et al. Recycling lithium-ion batteries from electric vehicles. Nature 575, 75--86 (2019).

\bibitem{ref3} Cucchiella, F., D’Adamo, I., Koh, S. C. L. \& Rosa, P. Recycling of WEEEs: An economic assessment of present and future e-waste streams. Renew. Sustain. Energy Rev. 51, 263--272 (2015).

\bibitem{ref4} Zeng, X., Mathews, J. A. \& Li, J. Urban mining of e-waste is becoming more cost-effective than virgin mining. Environ. Sci. Technol. 52, 4835--4841 (2018).

\bibitem{ref5} Tanskanen, P. Management and recycling of electronic waste. Acta Mater. 61, 1001--1011 (2013).

\bibitem{ref6} Wegener, K., Chen, W. H., Dietrich, F., Dröder, K. \& Kara, S. Robot-assisted disassembly for the recycling of electric vehicle batteries. Procedia CIRP 29, 716--721 (2015).

\bibitem{ref8} Driess, D. et al. PaLM-E: An embodied multimodal language model. in Proc. ICML 8469--8488 (2023).

\bibitem{ref9} Zitkovich, B., et al. RT-2: Vision-language-action models transfer web knowledge to robotic control. in Proc. CoRL, PMLR 229, 2165--2183 (2023).

\bibitem{ref10} Ahn, M. et al. Do as I can, not as I say: Grounding language in robotic affordances. Preprint at arXiv:2204.01691 (2022).

\bibitem{ref11} Liang, J. et al. Code as policies: Language model programs for embodied control. Preprint at arXiv:2209.07753 (2022).

\bibitem{ref12} Wang, G. et al. Voyager: An open-ended embodied agent with large language models. Preprint at arXiv:2305.16291 (2023).

\bibitem{ref13} Vemprala, S., Bonatti, R., Bucker, A. \& Kapoor, A. ChatGPT for robotics: Design principles and model abilities. Preprint at \href{https://arxiv.org/abs/2306.17582}{arXiv:2306.17582} (2023).

\bibitem{ref39} Rrapi, F., Portelli, B., Serra, G. \& Scalera, L. From AI foundations to Large Language Models: A survey on challenges and opportunities in collaborative robotics. Robot. Comput.-Integr. Manuf. 100, 103269 (2026). \url{https://doi.org/10.1016/j.rcim.2026.103269}.

\bibitem{ref16} Achiam, J., Held, D., Tamar, A. \& Abbeel, P. Constrained policy optimization. in Proc. ICML 22--31 (2017).

\bibitem{ref17} Altman, E. Constrained Markov Decision Processes (Chapman and Hall/CRC, 1999).

\bibitem{ref18} Ray, A., Achiam, J. \& Amodei, D. Benchmarking safe exploration in deep reinforcement learning. OpenAI Tech. Rep. (2019).

\bibitem{ref19} Tessler, C., Mankowitz, D. J. \& Mannor, S. Reward constrained policy optimization. in Proc. ICLR (2019).

\bibitem{ref20} Stooke, A., Achiam, J. \& Abbeel, P. Responsive safety in reinforcement learning by PID Lagrangian methods. in Proc. ICML 9133--9143 (2020).

\bibitem{ref21} Dalal, G., Dvijotham, K., Vecerik, M., Hester, T., Paduraru, C. \& Tassa, Y. Safe exploration in continuous action spaces. Preprint at \href{https://arxiv.org/abs/1801.08757}{arXiv:1801.08757} (2018).

\bibitem{ref22} Thananjeyan, B. et al. Recovery RL: Safe reinforcement learning with learned recovery zones. IEEE Robot. Autom. Lett. 6, 4915--4922 (2021).

\bibitem{ref23} Ames, A. D., Xu, X., Grizzle, J. W. \& Tabuada, P. Control barrier function based quadratic programs for safety critical systems. IEEE Trans. Autom. Control 62, 3861--3876 (2017).

\bibitem{ref24} Ames, A. D. et al. Control barrier functions: Theory and applications. in Proc. ECC 3420--3431 (2019).

\bibitem{ref40} Xiao, W. \& Belta, C. High-order control barrier functions. IEEE Trans. Autom. Control 67, 3655--3662 (2022). \url{https://doi.org/10.1109/TAC.2021.3105491}.

\bibitem{ref41} Lindemann, L. \& Dimarogonas, D. V. Control barrier functions for signal temporal logic tasks. IEEE Control Syst. Lett. 3, 96--101 (2019). \url{https://doi.org/10.1109/LCSYS.2018.2853182}.

\bibitem{ref33} Coumans, E. \& Bai, Y. PyBullet, a Python module for physics simulation for games, robotics and machine learning. \url{https://pybullet.org} (2016--2021).

\bibitem{ref31} Schulman, J., Wolski, F., Dhariwal, P., Radford, A. \& Klimov, O. Proximal policy optimization algorithms. Preprint at arXiv:1707.06347 (2017).

\bibitem{ref32} Haarnoja, T., Zhou, A., Abbeel, P. \& Levine, S. Soft actor-critic: Off-policy maximum entropy deep reinforcement learning with a stochastic actor. in Proc. ICML 1861--1870 (2018).

\bibitem{ref36} Anthropic. Model IDs and versioning. \url{https://platform.claude.com/docs/en/about-claude/models/model-ids-and-versions} (accessed 21 September 2026).

\bibitem{ref37} OpenAI. GPT-4o model: snapshots and aliases. \url{https://developers.openai.com/api/docs/models/gpt-4o} (accessed 21 September 2026).

\bibitem{ref38} Google. Models: Gemini API. \url{https://ai.google.dev/gemini-api/docs/models} (accessed 21 September 2026).
\end{thebibliography}
\end{document}
